\documentstyle[11pt]{article}
\textwidth 155mm
\textheight 220mm
\topmargin -24pt
\oddsidemargin 0.5cm
\evensidemargin 0.5cm
\def\Gr{Gr\"obner }
\begin{document}
\title{\bf Application of Involutive Division Methods to Polynomial
 and Differential Equations}
\author{Vladimir P. Gerdt \\
Laboratory of Computing Techniques and Automation\\
       Joint Institute for Nuclear Research\\
       141980 Dubna, Russia \\
       gerdt@jinr.dubna.su}
\date{}
\maketitle
\begin{abstract} A new algorithmic technique is presented
for constructing of polynomial and differential \Gr bases of the
special form called involutive by analogue with involutive systems of
partial differential equations. The corner-stone of the new technique is a
concept of involutive monomial division which provides the partition
of variables into multiplicative and non-multiplicative ones. This
partition leads to a self-consistent procedure for computation of an
involutive basis based on non-multiplicative prolongations and
multiplicative reductions. Every involutive division generates a
particular form of the computational procedure. Along with three
involutive divisions used by Thomas, Janet and Pommaret in analysis
of partial differential equations we consider some other divisions
of practical interest. It is shown that, given an involutive division
and an admissible monomial ordering, a minimal involutive basis is
uniquely defined for polynomial and linear partial differential equations
much as a reduced \Gr basis.
\end{abstract}

\end{document}
